\section{Discussion and limitations}\label{sec:discussion}

We presented three class specific conformal envelopes for classification with several predictors: a collapsed, a radial 
and a strip envelope. They share the same two stage calibration procedure and the same coverage statement, they use only the 
observed scores of a sample, and they are implemented in an open source Python package. In our experiments the coverage 
of all three constructions was close to the nominal level, and none of them was better than the others in every experiment.

The coverage statement of Section~\ref{sec:setup} has limits. It holds for each class and not conditionally on the input 
$x$, and it assumes that the samples of a class are exchangeable and that the missingness of the scores has the same 
distribution in the calibration and in the test data. Filtering the candidate classes before the prediction can make it 
fail whenever the filter removes the true class. If a class has few calibration samples, its threshold is infinite and 
its envelope accepts every sample, which gives large prediction sets. 

The experiments also have limits. The simulated data come from one generator in which scores are hidden completely at 
random, and the football data have no missing scores, so we did not study scores that are missing for a reason related 
to the class. We compared the three constructions only with each other, with the collapsed envelope as the simplest 
baseline. We did not compare them with other conformal methods that combine several predictors, in particular the method 
of \citet{ochoarivera2024conformal}. The parameters of the radial and strip envelopes were set to defaults and were not 
tuned. Comparisons with these methods, scores that are not missing at random, and tuning of the parameters are directions 
for future work.